% =====================================================================
%  Replacement narrative for Section 5.1 (continuous variables).
%
%  NOTE BEFORE USE: BET is absent from every number below because its
%  adapter is returning a quantity anti-correlated with dependence on the
%  symmetric schemes -- 37 per cent of its cells fall below a rejection
%  rate of 0.03 at a nominal 0.05, and on the circle its power RISES from
%  0.003 to 0.012 as the dependence weakens.  Once the adapter is fixed
%  and BET re-run, rebuild the tables and figures and reinstate it here.
%  The associate editor named BET explicitly, so it cannot be left out of
%  the submitted version.
% =====================================================================

We compared $B_n$ with ten competitors on thirteen continuous dependence
schemes, each indexed by a noise standard deviation ranging from $0{\cdot}1$
to $2{\cdot}0$ in steps of $0{\cdot}1$, at $n = 50$, $100$ and $200$, with
orthogonal polynomials, orders from rule~(12) and $10^4$ replications per
configuration. Ten of the schemes are those of the earlier literature
\citep{New09, Hel13, Sim14, Rud17}; three are added here. The \textit{lattice}
is a balanced $4\times4$ pattern with exactly uniform margins in which every
rank correlation vanishes identically, so that dependence is visible only once
the unit square is resolved; \textit{sine 8pi} quadruples the frequency of the
sine scheme, and is included because a low-order smooth span is blind to
high-frequency departures by construction and the design should contain the
case the method is worst at; and \textit{Clayton} supplies the lower-tail
dependence that none of the ten possessed. All critical values are
size-adjusted, taken from $2\times10^4$ null datasets per $(n,p,q)$ for the
statistics that depend on the data only through the ranks, and from
$5\times10^3$ permuted datasets per configuration for distance covariance and
HHG, which do not.

Table~\ref{tab:power_summary} reports the mean over the $260$ configurations at
each sample size and Fig.~\ref{fig:power} the full profiles. At $n = 50$, $B_n$
has the highest mean power of any test considered, $0{\cdot}462$ against
$0{\cdot}434$ for HHG and $0{\cdot}429$ for its atoms-based variant; at
$n = 100$ it is within $0{\cdot}033$ of the leader, and at $n = 200$ HHG leads
by $0{\cdot}054$. The ordering among the remaining competitors is stable, with
distance covariance, Hoeffding's statistic and the Cram\'er--von Mises
statistic of \citet{Gen19} within $0{\cdot}03$ of one another throughout, and
the maximal information coefficient and Chatterjee's $\xi$ trailing by a
further margin that widens with $n$.

What the averages conceal is that the tests are not competing on a common
axis, and the scheme-by-scheme profile is the more informative comparison.
$B_n$ is the most powerful test on every scheme in which the dependence is
neither monotone nor of high frequency. On the circle at $n = 100$ it rejects
with mean probability $0{\cdot}556$ against $0{\cdot}540$ for the nearest
competitor and $0{\cdot}078$ for distance covariance; on the lattice,
$0{\cdot}797$ against $0{\cdot}239$, $0{\cdot}264$ and $0{\cdot}216$ for
distance covariance, Hoeffding and \citet{Gen19} respectively, those three
being close to powerless by construction, since the lattice has no rank
correlation to detect. At $n = 50$ it leads on five of the thirteen schemes
and by margins of $0{\cdot}10$ to $0{\cdot}13$ on the structured ones.

The cases in which $B_n$ trails are equally systematic, and
Proposition~\ref{prop:lp_fixed} accounts for both of them exactly. The first is
monotone dependence: on the linear, Clayton, fourth-root and two-curves
schemes the departure is essentially one-dimensional, so the whole
noncentrality sits in the single coefficient $\theta_{11}$ while the test
spends $pq$ degrees of freedom and pays a critical value of order $pq$ for
them. On the linear scheme at $n = 200$, distance covariance attains
$0{\cdot}850$ against $0{\cdot}625$ for $B_n$. The remedy is in the theory
rather than outside it: with $p = q = 1$ the statistic is a monotone function
of Spearman's correlation and the deficit disappears, so what the table
measures is the price of not knowing in advance that the relation is monotone.
The second is high frequency: on the sine schemes the direction of departure
lies largely outside the span at $p = q = 3$, the captured energy
$\|\Pi_{pq}g\|^2$ is small, and the noncentrality with it. On \textit{sine 8pi}
at $n = 200$, $B_n$ reaches $0{\cdot}231$ against $0{\cdot}494$ for the
maximal information coefficient. Here too the span is the instrument: raising
the orders restores the power, at the cost of degrees of freedom elsewhere,
which is the trade-off \S\ref{sec:localpower} describes and
\S\,S4.4 of the Supplementary Material quantifies.

Two further observations bear on how the table should be read. First, the
relative standing of $B_n$ is strongest at the smallest sample size and
weakens as $n$ grows, because rule~(12) lets the span grow only as $n^{0{\cdot}31}$,
giving $p = q = 4$ at $n = 200$; the comparison at $n = 200$ is therefore
between a competitor using the whole sample and a test deliberately restricted
to sixteen directions. Second, of the eleven statistics only $B_n$ and $P_n$
carry a reference distribution in closed form. Distance covariance and HHG are
not distribution-free at all and were calibrated here by permutation within
every configuration, a cost borne by this study rather than by the tests as
reported; in an application that calibration is the user's to pay, and
\S\ref{sec:computation} shows what it amounts to. At $n = 10^4$ a single
evaluation of the HHG statistic takes $123$ seconds against $0{\cdot}007$ for
$B_n$, before any permutation at all. The comparison on power is close; the
comparison on what it costs to obtain a $p$-value is not.
